% GENERATED FILE. Edit canonical lessons, then run npm run generate:books.
% canonical-source-hash: fnv1a64:f20f6d2ad259c491
% canonical-lessons: PA-C123-lasan, PA-C123-adrak, PA-C123-narial, PA-C123-tarbuz, PA-C123-papita

\chapter{Garlic, Ginger, Coconut, Watermelon, Papaya}
\label{ch:pa-garlic-ginger-coconut-watermelon-papaya}
\hlchaptermodality{123}

\begin{chapteropening}
\textbf{By the end of this chapter:} \emph{I can say garlic, ginger, a coconut, a watermelon and a papaya.}

Complete the last lesson of chapter 123: i can say garlic, ginger, a coconut, a watermelon and a papaya.
\end{chapteropening}

\section[\texorpdfstring{\pa{ਲਸਣ} (lasaṇ)}{ਲਸਣ (lasaṇ)}]{\pa{ਲਸਣ} (lasaṇ) — garlic}

\label{lesson:PA-C123-lasan}



\begin{quote}
Before the new one: say the Punjabi for an aubergine, then the Punjabi for a cauliflower.
\end{quote}

\subsection*{You'll want to know: \pa{ਲਸਣ}}
\textbf{\pa{ਲਸਣ}} — \emph{lasaṇ} — \textquotedblleft{}garlic\textquotedblright{}.

\begin{grammarlens}[title={What you've built}]
One more word for this chapter.
\end{grammarlens}

\subsection*{Your turn}
\begin{itemize}
\raggedright
  \item \emph{Say it:} \emph{lasaṇ}
  \item \emph{Say it:} \emph{lasaṇ}, once more
  \item \emph{Say it:} say \emph{lasaṇ}
  \item \emph{Recall:} say the Punjabi for an aubergine, then the Punjabi for a cauliflower, then say \emph{lasaṇ} again
\end{itemize}

\begin{tcolorbox}[breakable,colback=teal!4,colframe=teal!35!black,title={Before you move on}]
What does \pa{ਲਸਣ} mean? (\textquotedblleft{}Garlic\textquotedblright{}.) Say it once more.
\end{tcolorbox}
\section[\texorpdfstring{\pa{ਅਦਰਕ} (adrak)}{ਅਦਰਕ (adrak)}]{\pa{ਅਦਰਕ} (adrak) — ginger}

\label{lesson:PA-C123-adrak}



\begin{quote}
Before the new one: say the Punjabi for a cauliflower, then the Punjabi for garlic.
\end{quote}

\subsection*{You'll want to know: \pa{ਅਦਰਕ}}
\textbf{\pa{ਅਦਰਕ}} — \emph{adrak} — \textquotedblleft{}ginger\textquotedblright{}.

\begin{grammarlens}[title={What you've built}]
One more word for this chapter.
\end{grammarlens}

\subsection*{Your turn}
\begin{itemize}
\raggedright
  \item \emph{Say it:} \emph{adrak}
  \item \emph{Say it:} \emph{adrak}, once more
  \item \emph{Say it:} say \emph{adrak}
  \item \emph{Recall:} say the Punjabi for a cauliflower, then the Punjabi for garlic, then say \emph{adrak} again
\end{itemize}

\begin{tcolorbox}[breakable,colback=teal!4,colframe=teal!35!black,title={Before you move on}]
What does \pa{ਅਦਰਕ} mean? (\textquotedblleft{}Ginger\textquotedblright{}.) Say it once more.
\end{tcolorbox}
\section[\texorpdfstring{\pa{ਨਾਰੀਅਲ} (nārīal)}{ਨਾਰੀਅਲ (nārīal)}]{\pa{ਨਾਰੀਅਲ} (nārīal) — a coconut}

\label{lesson:PA-C123-narial}



\begin{quote}
Before the new one: say the Punjabi for garlic, then the Punjabi for ginger.

In \emph{jo … oh …}, which half describes, and which points at the one meant? (\emph{Jo} describes; \emph{oh} points.)
\end{quote}

\subsection*{You'll want to know: \pa{ਨਾਰੀਅਲ}}
\textbf{\pa{ਨਾਰੀਅਲ}} — \emph{nārīal} — \textquotedblleft{}a coconut\textquotedblright{}.

\begin{grammarlens}[title={What you've built}]
One more word for this chapter.
\end{grammarlens}

\subsection*{Your turn}
\begin{itemize}
\raggedright
  \item \emph{Say it:} \emph{nārīal}
  \item \emph{Say it:} \emph{nārīal}, once more
  \item \emph{Say it:} say \emph{nārīal}
  \item \emph{Recall:} say the Punjabi for garlic, then the Punjabi for ginger, then say \emph{nārīal} again
\end{itemize}

\begin{tcolorbox}[breakable,colback=teal!4,colframe=teal!35!black,title={Before you move on}]
What does \pa{ਨਾਰੀਅਲ} mean? (\textquotedblleft{}A coconut\textquotedblright{}.) Say it once more.
\end{tcolorbox}
\section[\texorpdfstring{\pa{ਤਰਬੂਜ਼} (tarbūz)}{ਤਰਬੂਜ਼ (tarbūz)}]{\pa{ਤਰਬੂਜ਼} (tarbūz) — a watermelon}

\label{lesson:PA-C123-tarbuz}



\begin{quote}
Before the new one: say the Punjabi for ginger, then the Punjabi for a coconut.
\end{quote}

\subsection*{You'll want to know: \pa{ਤਰਬੂਜ਼}}
\textbf{\pa{ਤਰਬੂਜ਼}} — \emph{tarbūz} — \textquotedblleft{}a watermelon\textquotedblright{}.

\begin{grammarlens}[title={What you've built}]
One more word for this chapter.
\end{grammarlens}

\subsection*{Your turn}
\begin{itemize}
\raggedright
  \item \emph{Say it:} \emph{tarbūz}
  \item \emph{Say it:} \emph{tarbūz}, once more
  \item \emph{Say it:} say \emph{tarbūz}
  \item \emph{Recall:} say the Punjabi for ginger, then the Punjabi for a coconut, then say \emph{tarbūz} again
\end{itemize}

\begin{tcolorbox}[breakable,colback=teal!4,colframe=teal!35!black,title={Before you move on}]
What does \pa{ਤਰਬੂਜ਼} mean? (\textquotedblleft{}A watermelon\textquotedblright{}.) Say it once more.
\end{tcolorbox}
\section[\texorpdfstring{\pa{ਪਪੀਤਾ} (papītā)}{ਪਪੀਤਾ (papītā)}]{\pa{ਪਪੀਤਾ} (papītā) — a papaya}

\label{lesson:PA-C123-papita}



\begin{quote}
Before the new one: say the Punjabi for a coconut, then the Punjabi for a watermelon.
\end{quote}

\subsection*{You'll want to know: \pa{ਪਪੀਤਾ}}
\textbf{\pa{ਪਪੀਤਾ}} — \emph{papītā} — \textquotedblleft{}a papaya\textquotedblright{}.

\begin{grammarlens}[title={What you've built}]
One more word for this chapter.
\end{grammarlens}

\subsection*{Your turn}
\begin{itemize}
\raggedright
  \item \emph{Say it:} \emph{papītā}
  \item \emph{Say it:} \emph{papītā}, once more
  \item \emph{Say it:} say \emph{papītā}
  \item \emph{Recall:} say the Punjabi for a coconut, then the Punjabi for a watermelon, then say \emph{papītā} again
\end{itemize}

\begin{tcolorbox}[breakable,colback=teal!4,colframe=teal!35!black,title={Before you move on}]
What does \pa{ਪਪੀਤਾ} mean? (\textquotedblleft{}A papaya\textquotedblright{}.) Say it once more.
\end{tcolorbox}
